% Generated draft: requires mathematical and source review.
\documentclass[11pt]{article}
\usepackage[margin=25mm]{geometry}
\usepackage{fontspec}
\setmainfont{Latin Modern Roman}
\usepackage{amsmath,amssymb,mathtools,unicode-math}
\setmathfont{Latin Modern Math}
\usepackage{enumitem,needspace,xurl,hyperref}
\setlength{\parindent}{0pt}
\setlength{\parskip}{5pt}
\setlength{\emergencystretch}{3em}
\newcommand{\N}{\mathbb{N}}
\newcommand{\Z}{\mathbb{Z}}
\newcommand{\Q}{\mathbb{Q}}
\newcommand{\R}{\mathbb{R}}
\newcommand{\C}{\mathbb{C}}
\newcommand{\F}{\mathbb{F}}
\newcommand{\A}{\mathbb{A}}
\newcommand{\PP}{\mathbb{P}}
\newcommand{\E}{\mathbb{E}}
\newcommand{\eps}{\varepsilon}
\newcommand{\dd}{\,\mathrm{d}}
\begin{document}
\section*{UnsolvedMath ID: 10400202}
\textbf{Source problem code:} catalogue entry
\subsection{Definitions and mathematical statement}
\emph{Problem.} Construct the LMO invariant, or the theory of finite-type invariants, in terms of arrow diagrams. The source does not specify the arrow diagrams, their skeletons or labels, their relations, or the coefficient and completion conventions. These choices, together with the construction and its invariance, form part of the requested work.

For the LMO alternative, an arrow-diagram construction must recover the LMO invariant in the source's intended manifold class and normalization. For example, if it produces a value \(Z_{\mathrm{arr}}(M)\) in a proposed diagram space \(\mathcal A_{\mathrm{arr}}\), a suitable comparison map \(\Phi\) would have to identify it with the usual invariant:
\[
 \Phi\bigl(Z_{\mathrm{arr}}(M)\bigr)=Z_{\mathrm{LMO}}(M).
\]
This is a way to express the recovery requirement, not a specified source diagram algebra or an assertion that such a map exists. The alternative request is an arrow-diagram description of the finite-type theory, whose filtration, degree conventions, and class of manifolds must likewise be specified. The word ``or'' is retained: the excerpt does not demand both constructions simultaneously.

The supplied record continues into the separate section ``11.5 Refinements of the LMO invariant (T. Le).'' That context refers to the remark in Problem 11.1, says the LMO invariant is weak when \(b_1(M)>0\), and states
\[
 Z_{\mathrm{LMO}}(M)=1\quad\text{when }b_1(M)>3,
 \qquad b_1(M)=\dim_\Q H_1(M;\Q).
\]
The formula is retained as a source claim with its unprovided normalization and manifold conventions; it is not an independently verified theorem here. The following two, unprovided problems are said to seek stronger refinements, especially for positive first Betti number. This adjacent discussion does not add a stronger-than-LMO requirement to Problem 11.6, and its T. Le attribution is not assigned to the numbered arrow-diagram problem without evidence.
\subsection{Short English statement}
Give an arrow-diagram construction of the LMO invariant or of the finite-type invariant theory, including the relations and comparison needed to recover the intended theory. The following paragraph about positive-Betti-number refinements is preserved as separate context.
\subsection{Sources}
\begin{itemize}
\item Local UnsolvedMath registry, permanent ID 10400202, source code catalogue entry, Problem 11.6. No individual proposer is supplied for this numbered problem.
\item Ohtsuki, \emph{Problems on invariants of knots and 3-manifolds} (2002), Problem 11.6, PDF page 156, followed in the supplied extraction by section 11.5, ``Refinements of the LMO invariant (T. Le),'' referring to Problem 11.1. Supplied URL: \url{https://msp.org/gtm/2002/04/gtm-2002-04-024s.pdf}; recorded access 2026-07-29; extraction pdf-font-structure. The collection distinguishes 231 numbered bold headings from references and surrounding exposition.
\item The registry labels the item open and OPEN-TRIAGE. Background status evidence remains NEEDS\_REVIEW pending a current primary-literature audit; redistribution terms of the publicly accessible author/publisher copy also require review.
\end{itemize}
\end{document}
